\vfill\pagebreak
\section{Putting it together: a guessing game}
\label{sec:control:guessing_game}

The program below picks a number between 1 and 100 and lets the user guess it,
answering each guess with ``too small'' or ``too big'' until the number is
found.

\begin{lstlisting}[style=coloredverbatim, caption=A guessing game, label=lst:control:guessing_game]
use std::io;
use std::rand;

fn main() {
  let secret = uniform(1, 100);
  let mut tries = 0;

  println("I am thinking of a number between 1 and 100.");
  loop {
    let guess = read!i32(ask-> "Your guess: ");
    tries += 1;

    if guess < secret {
      println("Too small.");
    } else if guess > secret {
      println("Too big.");
    } else {
      println("Found it in ", tries, " tries!");
      break;
    }
  }
}
\end{lstlisting}
\vspace{-10pt}%
\begin{lstlisting}[style=bashVerb, escapechar=@+]
@+\textcolor{teal!80}{alice@dev:\mtilde}@+$ gyc main.yr -o main
@+\textcolor{teal!80}{alice@dev:\mtilde}@+$ ./main
I am thinking of a number between 1 and 100.
Your guess: 50
Too small.
Your guess: 75
Too big.
Your guess: 62
Found it in 3 tries!
\end{lstlisting}

Read from the top, each construction has a job.

\begin{itemize}
\item \token{uniform} draws the secret on line~5, so it differs on every run.
  Every decision the program takes depends on it, or on what the user types.
\item \token{tries} is declared before the loop, so that it keeps counting from
  one guess to the next.
\item The \token{loop} on line~9 repeats the question. It is a \token{loop}
  rather than a \token{while}, because the program can only decide to stop after
  a guess has been read and compared.
\item The chain of \token{if}, \token{else if} and \token{else} sorts each
  guess into one of three cases. Exactly one of them runs.
\item The \token{break} on line~19, in the last case only, is the one way out
  of the loop.
\end{itemize}

\notebox{The session above guesses by halving: 50 is the middle of 1 to 100, 75
  the middle of what is left, and so on. Each answer halves the numbers still
  possible, and seven halvings take a hundred possibilities down to one, so a
  player who always halves never needs more than seven tries. The last exercise
  of the chapter builds on that.}
